\documentclass[10pt,a4paper]{amsart}
\usepackage[margin=19mm]{geometry}
\usepackage{amsmath,amssymb,mathtools}
\usepackage[scr=rsfs,cal=euler]{mathalfa}
\usepackage{xfrac}
\usepackage[unicode,hidelinks,bookmarks=false]{hyperref}
\newcommand{\ie}{i.e.\ }
\newcommand{\cf}{cf.\ }
\newcommand{\resp}{respectively\ }
\newcommand{\Subsection}{\subsection}
\providecommand{\nobreakdash}{\nobreak\hbox{-}}
\setlength{\emergencystretch}{2em}
\multlinegap0pt
\usepackage[shortlabels]{enumitem}
\usepackage{float}
\usepackage{csquotes}
\usepackage{amssymb}
\usepackage{xcolor}
\usepackage{tikz}
\newtheorem{theorem}{Theorem}
\newtheorem{lemma}{Lemma}
\newcommand{\Z}{\mathbb{Z}}
\title{Weyl Groups and the Modified Kostant Game}
\author{Alexander Caviedes Castro, Juan Sebasti\'an Cort\'es-Cruz}
\date{Source: arXiv:2605.11449v1 (CC BY 4.0)}
\begin{document}
\maketitle

\noindent\emph{Source and licence.} Weyl Groups and the Modified Kostant Game. Version arXiv:2605.11449v1 (CC BY 4.0), distributed by arXiv under the Creative Commons Attribution 4.0 International licence, http://creativecommons.org/licenses/by/4.0/, as recorded in the arXiv licence metadata for this exact version and shown on the abstract page https://arxiv.org/abs/2605.11449v1. Full licence notice: \texttt{licenses/arxiv-2605.11449-CC-BY-4.0.txt}.\par\smallskip

\noindent\textit{Editorial scope.} Counted unit: the article's theorem inequality (source0.tex lines 1352--1376). The article's complete original proof of it is delivered with it (source0.tex lines 1378--1498, one \texttt{proof} environment of 121 lines). No mathematics is written, repaired or re-derived: the statement and the proof are the article's own lines, in the article's own notation, and no retained line is substituted or re-flowed.  Retained uncounted as the source-local context the proof needs: lines 932--937; lines 1337--1342; lines 1344--1350.  Nothing was omitted from any retained span, so the package carries no omission record.  No retained line contains a citation, so the wrapper carries no bibliography and no bibliography entry.  No retained line contains a figure environment, an image-inclusion command or drawing code, so no image is delivered and the package declares no asset dependency.  Editorial formatting only, in addition: an AMS \texttt{amsart} preamble that restates the article's own declarations and macros used by the retained lines, namely the environments \texttt{theorem}, \texttt{lemma} on separate counters, together with the standard packages those lines need.  The retained environments print in this document as Theorem 1, Lemma 1 in order of appearance, whereas the article prints the same items with the numbering of its own full document.  The complete native source of the article as distributed in the arXiv e-print for this exact version is bundled unchanged as \texttt{sources/200357/source0.tex}.
\begin{theorem}[Root Counting Identity]\label{thm:root_counting_identity}
	Let $I \subseteq S$ be an arbitrary subset of modified vertices, and $J = S \setminus I$. The Modified Kostant Game terminates in a unique final configuration $\mathbf{c}_{final}^{(I)}$. Moreover, the total number of chips in this final state equals the sum of the $I$-heights of all positive coroots that are not in the subsystem generated by $J$:
	\begin{equation}
		|\mathbf{c}_{final}^{(I)}| = \sum_{\gamma^\vee \in (\Phi^\vee)^+ \setminus (\Phi_J^\vee)^+} \mathrm{ht}_I(\gamma^\vee).
	\end{equation}
\end{theorem}

\begin{lemma}\label{root sum lemma}
Let $\beta$ be a simple root. Then
\[
\sum_{\alpha\in \Phi^+}\langle \alpha, \beta^\vee \rangle=2.
\]
\end{lemma}

\begin{lemma}\label{lemma:string}
Let $\alpha$ and $\beta$ be non-proportional roots. Let $r, q\in
\Z_{\geq 0}$ be the largest integers for which $\beta-r\alpha\in \Phi,
\beta+q\alpha\in \Phi$. Then $\beta+i\alpha\in \Phi$ for all $-r\leq i\leq
q$. The set $\{\beta+i\alpha : i\in \Z\}$ is called the
\textbf{$\alpha$-string through $\beta$}.
\end{lemma}

\begin{theorem}
For any subset $\Delta_P \subset \Delta$, the inequality
\begin{equation}\label{strong inequality}
\sum_{\beta\in \Delta \smallsetminus \Delta_P}
\Biggl(
\sum_{\alpha\in \Phi^+\smallsetminus \Phi_{\Delta_P}^+}
\langle \alpha,\beta^\vee\rangle-1
\Biggr)
\leq |\Phi^+\smallsetminus \Phi_{\Delta_P}^+|
\end{equation}
holds. As a consequence,
\[
|\Delta \smallsetminus \Delta_P|\cdot
\Biggl(
\gcd_{\beta\in \Delta \smallsetminus \Delta_P}
\Biggl(
\sum_{\alpha\in \Phi^+\smallsetminus \Phi_{\Delta_P}^+}
\langle \alpha,\beta^\vee\rangle
\Biggr)
-1
\Biggr)
\le
|\Phi^+\smallsetminus \Phi_{\Delta_P}^+|.
\]
\end{theorem}

\begin{proof}
Consider the subset $\tilde{\Phi}_{P}\subset (\Phi^\vee)^+$ consisting of all positive coroots of the form
\[
\beta^\vee+\sum_{\alpha_j\in \Delta_P} n_j\alpha_j^\vee,
\]
where $\beta\in \Delta \smallsetminus \Delta_P$ and $n_j\in\mathbb Z_{\ge 0}$.
Clearly,
\[
\tilde{\Phi}_{P}\subset (\Phi^\vee)^+\smallsetminus (\Phi_{\Delta_P}^\vee)^+.
\]
Thus, it is enough to prove that
\[
\sum_{\beta\in \Delta \smallsetminus \Delta_P}
\Biggl(
\sum_{\alpha\in \Phi^+\smallsetminus \Phi_{\Delta_P}^+}
\langle \alpha,\beta^\vee\rangle-1
\Biggr)
\le |\tilde{\Phi}_{P}|.
\]

Fix $\beta \in \Delta \smallsetminus \Delta_P$. Consider the Dynkin diagram associated with $\Delta_P$, and let $\Gamma_{\Delta_P}$ be the corresponding subgraph. Denote by
\[
\Gamma_{P_j} = \Gamma_{P_j}(\beta)
\]
the connected components of $\Gamma_{\Delta_P}$ that are adjacent to $\beta$. Note that there are at most three such components. Suppose that for a component $\Gamma_{P_j}$ there are $k$ arrows pointing from the simple root in $\Gamma_{P_j}$ adjacent to $\beta$ towards $\beta$.

For each component $\Gamma_{P_j}$, consider a string of coroots
\[
\beta^\vee,\,
\beta^\vee+\gamma_1^j,\,
\ldots,\,
\beta^\vee+\gamma_{l_j}^j
\]
obtained by playing the modified Kostant game on $\Gamma_{P_j}\cup\{\beta\}$ with active set $I=\{\beta\}$ and with the arrow multiplicities dictated by the Dynkin diagram, until termination. Each $\gamma_i^j$ is a positive linear combination of coroots in $(\Phi_{P_j}^\vee)^+$.
If $\beta$ is not adjacent to any vertex of $\Delta_P$, we simply take the string consisting only of $\beta^\vee$.

The mechanics of the Kostant game imply that these strings can be concatenated together. For instance, the following is a string of coroots:
\[
\beta^\vee,\,
\beta^\vee+\gamma_1^1,\,
\ldots,\,
\beta^\vee+\gamma_{l_1}^1,\,
\beta^\vee+\gamma_{l_1}^1+\gamma_1^2,\,
\ldots,\,
\beta^\vee+\sum_j\gamma_{l_j}^j.
\]
By Lemma~\ref{lemma:string}, this string can be completed so that the resulting coroots have consecutive heights. Hence its total length is
\[
1+\sum_j l_j
=1+\sum_j l_j(\beta).
\]
Moreover, every root corresponding to a coroot in this string belongs to $\tilde{\Phi}_{P}$.

It remains to identify this quantity with
\[
\sum_{\alpha\in \Phi^+\smallsetminus \Phi_{\Delta_P}^+}
\langle \alpha,\beta^\vee\rangle-1.
\]
By Lemma~\ref{root sum lemma},
\begin{align*}
\sum_{\alpha\in \Phi^+\smallsetminus \Phi_{\Delta_P}^+}
\langle \alpha,\beta^\vee\rangle-1
&=
1-\sum_{\alpha\in \Phi_{\Delta_P}^+}\langle \alpha,\beta^\vee\rangle \\
&=
1-\sum_j
\sum_{\alpha\in \Phi_{P_j}^+}\langle \alpha,\beta^\vee\rangle.
\end{align*}

For each component $\Gamma_{P_j}$, write
\[
\sum_{\alpha\in \Phi_{P_j}^+}\alpha
=
h_j\alpha_j+\cdots,
\]
where $\alpha_j$ is the simple root in $\Gamma_{P_j}$ adjacent to $\beta$.
Since $\beta^\vee$ is orthogonal to all simple roots in $\Delta_P$ except $\alpha_j$, we obtain
\[
\sum_{\alpha\in \Phi_{P_j}^+}\langle \alpha,\beta^\vee\rangle
=
h_j\langle \alpha_j,\beta^\vee\rangle
=
-kh_j.
\]
On the other hand, by the root counting identity from Theorem~\ref{thm:root_counting_identity} applied with $I=\{\beta\}$, the height of the final coroot of the game on $\Gamma_{P_j}\cup\{\beta\}$ is
\[
\operatorname{ht}(\beta^\vee+\gamma_{l_j}^j)=1+kh_j.
\]
Therefore,
\[
-\sum_{\alpha\in \Phi_{P_j}^+}\langle \alpha,\beta^\vee\rangle
=
\operatorname{ht}(\beta^\vee+\gamma_{l_j}^j)-1.
\]
Substituting this into the previous expression gives
\begin{align*}
\sum_{\alpha\in \Phi^+\smallsetminus \Phi_{\Delta_P}^+}
\langle \alpha,\beta^\vee\rangle-1
&=
1+\sum_j
\bigl(\operatorname{ht}(\beta^\vee+\gamma_{l_j}^j)-1\bigr) \\
&=
\operatorname{ht}\Bigl(\beta^\vee+\sum_j\gamma_{l_j}^j\Bigr) \\
&=
1+\sum_j l_j(\beta).
\end{align*}

Summing over all $\beta\in \Delta \smallsetminus \Delta_P$, we conclude that
\[
\sum_{\beta\in \Delta \smallsetminus \Delta_P}
\Biggl(
\sum_{\alpha\in \Phi^+\smallsetminus \Phi_{\Delta_P}^+}
\langle \alpha,\beta^\vee\rangle-1
\Biggr)
=
\sum_{\beta\in \Delta \smallsetminus \Delta_P}
\bigl(1+\sum_j l_j(\beta)\bigr)
\le |\tilde{\Phi}_{P}|,
\]
which proves the theorem.
\end{proof}

\end{document}
